
\subsection{Element Instance Generation}

\scriptsize
\texttt{---Goal---\\
Given a text document that is potentially relevant to this activity and a list of entity types, identify all entities of those types from the text and all relationships among the identified entities.\\
\\
---Steps---\\
1. Identify all entities. For each identified entity, extract the following information:\\
- entity\_name: Name of the entity, capitalized\\
- entity\_type: One of the following types: [\{entity\_types\}]\\
- entity\_description: Comprehensive description of the entity's attributes and activities\\
\\
Format each entity as ("entity"\{tuple\_delimiter\}<entity\_name>\{tuple\_delimiter\}<entity\_type>\{tuple\_\\
delimiter\}<entity\_description>\\
\\
2. From the entities identified in step 1, identify all pairs of (source\_entity, target\_entity) that are *clearly related* to each other\\
For each pair of related entities, extract the following information:\\
- source\_entity: name of the source entity, as identified in step 1\\
- target\_entity: name of the target entity, as identified in step 1\\
- relationship\_description: explanation as to why you think the source entity and the target entity are related to each other\\
- relationship\_strength: a numeric score indicating strength of the relationship between the source entity and target entity\\
\\
Format each relationship as ("relationship"\{tuple\_delimiter\}<source\_entity>\{tuple\_delimiter\}<target\_\\
entity>\{tuple\_delimiter\}<relationship\_description>\{tuple\_delimiter\}<relationship\_strength>)\\
\\
3. Return output in English as a single list of all the entities and relationships identified in steps 1 and 2. Use **\{record\_delimiter\}** as the list delimiter.\\
\\
4. When finished, output \{completion\_delimiter\}\\
\\
---Examples---\\
\\
Entity\_types: ORGANIZATION,PERSON\\
\\
Input:\\
\\
The Fed is scheduled to meet on Tuesday and Wednesday, with the central bank planning to release its latest policy decision on Wednesday at 2:00 p.m. ET, followed by a press conference where Fed Chair Jerome Powell will take questions. Investors expect the Federal Open Market Committee to hold its benchmark interest rate steady in a range of 5.25\%-5.5\%.\\
\\
Output:\\
\\
("entity"\{tuple\_delimiter\}FED\{tuple\_delimiter\}ORGANIZATION\{tuple\_delimiter\}The Fed is the Federal Reserve, which is setting interest rates on Tuesday and Wednesday)\\
\{record\_delimiter\}\\
("entity"\{tuple\_delimiter\}JEROME POWELL\{tuple\_delimiter\}PERSON\{tuple\_delimiter\}Jerome Powell is the chair of the Federal Reserve)\\
\{record\_delimiter\}\\
("entity"\{tuple\_delimiter\}FEDERAL OPEN MARKET COMMITTEE\{tuple\_delimiter\}ORGANIZATION\{tuple\_delimiter\}The Federal Reserve committee makes key decisions about interest rates and the growth of the United States money supply)\\
\{record\_delimiter\}\\
("relationship"\{tuple\_delimiter\}JEROME POWELL\{tuple\_delimiter\}FED\{tuple\_delimiter\}Jerome Powell is the Chair of the Federal Reserve and will answer questions at a press conference\{tuple\_delimiter\}9)\\
\{completion\_delimiter\}\\
}
\textit{\\
...More examples...\\
\\}
\texttt{---Real Data---\\
\\
Entity\_types: \{entity\_types\}\\
Input:\\
\\
\{input\_text\}\\
\\
Output:}

\subsection{Community Summary Generation}

\texttt{---Role--- \\
You are an AI assistant that helps a human analyst to perform general information discovery. Information discovery is the process of identifying and assessing relevant information associated with certain entities (e.g., organizations and individuals) within a network.\\
\\
---Goal---\\
Write a comprehensive report of a community, given a list of entities that belong to the community as well as their relationships and optional associated claims. The report will be used to inform decision-makers about information associated with the community and their potential impact. The content of this report includes an overview of the community's key entities, their legal compliance, technical capabilities, reputation, and noteworthy claims.\\
\\
---Report Structure---\\
\\
The report should include the following sections:\\
\\
- TITLE: community's name that represents its key entities - title should be short but specific. When possible, include representative named entities in the title.\\
- SUMMARY: An executive summary of the community's overall structure, how its entities are related to each other, and significant information associated with its entities.\\
- IMPACT SEVERITY RATING: a float score between 0-10 that represents the severity of IMPACT posed by entities within the community.  IMPACT is the scored importance of a community.\\
- RATING EXPLANATION: Give a single sentence explanation of the IMPACT severity rating.\\
- DETAILED FINDINGS: A list of 5-10 key insights about the community. Each insight should have a short summary followed by multiple paragraphs of explanatory text grounded according to the grounding rules below. Be comprehensive.\\
\\
Return output as a well-formed JSON-formatted string with the following format:\\
    \{\{\\
        "title": <report\_title>,\\
        "summary": <executive\_summary>,\\
        "rating": <impact\_severity\_rating>,\\
        "rating\_explanation": <rating\_explanation>,\\
        "findings": [\\
            \{\{\\
                "summary":<insight\_1\_summary>,\\
                "explanation": <insight\_1\_explanation>\\
            \}\},\\
            \{\{\\
                "summary":<insight\_2\_summary>,\\
                "explanation": <insight\_2\_explanation>\\
            \}\}\\
        ]\\
    \}\}\\
\\
---Grounding Rules---\\
\\
Points supported by data should list their data references as follows:\\
\\
"This is an example sentence supported by multiple data references [Data: <dataset name> (record ids); <dataset name> (record ids)]."\\
\\
Do not list more than 5 record ids in a single reference. Instead, list the top 5 most relevant record ids and add "+more" to indicate that there are more.\\
\\
For example:\\
"Person X is the owner of Company Y and subject to many allegations of wrongdoing [Data: Reports (1), Entities (5, 7); Relationships (23); Claims (7, 2, 34, 64, 46, +more)]."\\
\\
where 1, 5, 7, 23, 2, 34, 46, and 64 represent the id (not the index) of the relevant data record.\\
\\
Do not include information where the supporting evidence for it is not provided.\\
\\
---Example---\\
\\
Input:\\
\\
Entities\\
\\
id,entity,description\\
5,VERDANT OASIS PLAZA,Verdant Oasis Plaza is the location of the Unity March\\
6,HARMONY ASSEMBLY,Harmony Assembly is an organization that is holding a march at Verdant Oasis Plaza\\
\\
Relationships\\
\\
id,source,target,description\\
37,VERDANT OASIS PLAZA,UNITY MARCH,Verdant Oasis Plaza is the location of the Unity March\\
38,VERDANT OASIS PLAZA,HARMONY ASSEMBLY,Harmony Assembly is holding a march at Verdant Oasis Plaza\\
39,VERDANT OASIS PLAZA,UNITY MARCH,The Unity March is taking place at Verdant Oasis Plaza\\
40,VERDANT OASIS PLAZA,TRIBUNE SPOTLIGHT,Tribune Spotlight is reporting on the Unity march taking place at Verdant Oasis Plaza\\
41,VERDANT OASIS PLAZA,BAILEY ASADI,Bailey Asadi is speaking at Verdant Oasis Plaza about the march\\
43,HARMONY ASSEMBLY,UNITY MARCH,Harmony Assembly is organizing the Unity March\\
\\
Output:\\
\{\{\\
    "title": "Verdant Oasis Plaza and Unity March",\\
    "summary": "The community revolves around the Verdant Oasis Plaza, which is the location of the Unity March. The plaza has relationships with the Harmony Assembly, Unity March, and Tribune Spotlight, all of which are associated with the march event.",\\
    "rating": 5.0,\\
    "rating\_explanation": "The impact severity rating is moderate due to the potential for unrest or conflict during the Unity March.",\\
    "findings": [\\
        \{\{\\
            "summary": "Verdant Oasis Plaza as the central location",\\
            "explanation": "Verdant Oasis Plaza is the central entity in this community, serving as the location for the Unity March. This plaza is the common link between all other entities, suggesting its significance in the community. The plaza's association with the march could potentially lead to issues such as public disorder or conflict, depending on the nature of the march and the reactions it provokes. [Data: Entities (5), Relationships (37, 38, 39, 40, 41,+more)]"\\
        \}\},\\
        \{\{\\
            "summary": "Harmony Assembly's role in the community",\\
            "explanation": "Harmony Assembly is another key entity in this community, being the organizer of the march at Verdant Oasis Plaza. The nature of Harmony Assembly and its march could be a potential source of threat, depending on their objectives and the reactions they provoke. The relationship between Harmony Assembly and the plaza is crucial in understanding the dynamics of this community. [Data: Entities(6), Relationships (38, 43)]"\\
        \}\},\\
        \{\{\\
            "summary": "Unity March as a significant event",\\
            "explanation": "The Unity March is a significant event taking place at Verdant Oasis Plaza. This event is a key factor in the community's dynamics and could be a potential source of threat, depending on the nature of the march and the reactions it provokes. The relationship between the march and the plaza is crucial in understanding the dynamics of this community. [Data: Relationships (39)]"\\
        \}\},\\
        \{\{\\
            "summary": "Role of Tribune Spotlight",
            "explanation": "Tribune Spotlight is reporting on the Unity March taking place in Verdant Oasis Plaza. This suggests that the event has attracted media attention, which could amplify its impact on the community. The role of Tribune Spotlight could be significant in shaping public perception of the event and the entities involved. [Data: Relationships (40)]"\\
        \}\}\\
    ]\\
\}\}\\
\\
---Real Data---\\
\\
Use the following text for your answer. Do not make anything up in your answer.\\
\\
Input:\\
\\
\{input\_text\}\\}
\textit{\\
...Report Structure and Grounding Rules Repeated...\\
\\}
\texttt{Output:}

\subsection{Community Answer Generation}

\texttt{---Role---\\
\\
You are a helpful assistant responding to questions about a dataset by synthesizing perspectives from multiple analysts.\\
\\
---Goal---\\
\\
Generate a response of the target length and format that responds to the user's question, summarize all the reports from multiple analysts who focused on different parts of the dataset, and incorporate any relevant general knowledge.\\
\\
Note that the analysts' reports provided below are ranked in the **descending order of helpfulness**.\\
\\
If you don't know the answer, just say so. Do not make anything up.\\
\\
The final response should remove all irrelevant information from the analysts' reports and merge the cleaned information into a comprehensive answer that provides explanations of all the key points and implications appropriate for the response length and format.\\
\\
Add sections and commentary to the response as appropriate for the length and format. Style the response in markdown.\\
\\
The response shall preserve the original meaning and use of modal verbs such as "shall", "may" or "will".\\
\\
The response should also preserve all the data references previously included in the analysts' reports, but do not mention the roles of multiple analysts in the analysis process.\\
\\
Do not list more than 5 record ids in a single reference. Instead, list the top 5 most relevant record ids and add "+more" to indicate that there are more.\\
\\
For example:\\
\\
"Person X is the owner of Company Y and subject to many allegations of wrongdoing [Data: Reports (2, 7, 34, 46, 64, +more)]. He is also CEO of company X [Data: Reports (1, 3)]"\\
\\
where 1, 2, 3, 7, 34, 46, and 64 represent the id (not the index) of the relevant data record.\\
\\
Do not include information where the supporting evidence for it is not provided.\\
\\
---Target response length and format---\\
\\
\{response\_type\}\\
\\
---Analyst Reports---\\
\\
\{report\_data\}\\
\\}
\textit{\\
...Goal and Target response length and format repeated...\\
\\}
\texttt{Add sections and commentary to the response as appropriate for the length and format. Style the response in markdown.\\
\\
Output:}

\subsection{Global Answer Generation}

\texttt{---Role---\\
\\
You are a helpful assistant responding to questions about data in the tables provided.\\
\\
---Goal---\\
\\
Generate a response of the target length and format that responds to the user's question, summarize all relevant information in the input data tables appropriate for the response length and format, and incorporate any relevant general knowledge.\\
\\
If you don't know the answer, just say so. Do not make anything up.\\
\\
The response shall preserve the original meaning and use of modal verbs such as "shall", "may" or "will".\\
\\
Points supported by data should list the relevant reports as references as follows:\\
\\
"This is an example sentence supported by data references [Data: Reports (report ids)]"\\
\\}
\textit{Note: the prompts for SS (semantic search) and TS (text summarization) conditions use "Sources" in place of "Reports" above.\\}
\texttt{\\
Do not list more than 5 record ids in a single reference. Instead, list the top 5 most relevant record ids and add "+more" to indicate that there are more.\\
\\
For example:\\
\\
"Person X is the owner of Company Y and subject to many allegations of wrongdoing [Data: Reports (2, 7, 64, 46, 34, +more)]. He is also CEO of company X [Data: Reports (1, 3)]"\\
\\
where 1, 2, 3, 7, 34, 46, and 64 represent the id (not the index) of the relevant data report in the provided tables.\\
\\
Do not include information where the supporting evidence for it is not provided.\\
\\
At the beginning of your response, generate an integer score between 0-100 that indicates how **helpful** is this response in answering the user's question. Return the score in this format: <ANSWER\_HELPFULNESS> score\_value </ANSWER\_HELPFULNESS>.\\
\\
---Target response length and format---\\
\\
\{response\_type\}\\
\\
---Data tables---\\
\\
\{context\_data\}\\}
\textit{\\
...Goal and Target response length and format repeated...\\
\\}
Output: